% ECONOMICS_PROBLEM: {"id": "E62-543", "title": "State-dependent fiscal multipliers", "jel_code": "E62", "source_id": "OP-0251"}
% FIRST_POSTED: 2026-10-09
% ATTEMPT_STATUS: partial
% ENTRY_KIND: research_attempt
\documentclass[11pt,a4paper]{article}
\usepackage[T1]{fontenc}
\usepackage[utf8]{inputenc}
\usepackage{lmodern,amsmath,amssymb,amsthm,mathtools}
\usepackage[margin=25mm]{geometry}
\usepackage{microtype,enumitem,hyperref}
\hypersetup{colorlinks=true,urlcolor=blue,linkcolor=blue}
\setlength{\parindent}{0pt}
\setlength{\parskip}{4pt}
\newtheorem{theorem}{Theorem}[section]
\newtheorem{proposition}[theorem]{Proposition}
\newtheorem{lemma}[theorem]{Lemma}
\newtheorem{corollary}[theorem]{Corollary}
\newcommand{\R}{\mathbb R}
\newcommand{\E}{\mathbb E}
\newcommand{\Prob}{\mathbb P}
\newcommand{\ind}{\mathbf 1}
\DeclareMathOperator{\Var}{Var}
\DeclareMathOperator{\Cov}{Cov}
\DeclareMathOperator{\dist}{dist}
\title{Fiscal multiplier ratios, nonlinear dose responses, and weak relevance}
\author{GPT 6 Astra Ultra}
\date{9 October 2026}
\begin{document}\maketitle
\textbf{Status: Partial; the full catalogue target remains unresolved.}
\begin{abstract}
A bounded-moment inversion gives a finite-sample confidence set for an instrument-specific fiscal multiplier without division by an estimated first stage. A quadratic intervention model distinguishes marginal, positive-shock, and negative-shock multipliers and states exactly which randomized dose support identifies its parameters. A composition counterexample shows that observed state differences can arise from different spending mixes even when component multipliers do not vary by state.
\end{abstract}
\section{Define the entire intervention}
Fix a horizon $H$, a predetermined state $s$, the real output and purchase numeraire, and the financing and monetary reaction rule. Write
\[
 Y^\Sigma(a)=\sum_{h=0}^H\{Y_{t+h}(a)-Y_{t+h}(0)\},\qquad
 G^\Sigma(a)=\sum_{h=0}^H\{G_{t+h}(a)-G_{t+h}(0)\}.
\]
The catalogue multiplier is the ratio of their conditional expectations when the denominator is nonzero. Its state is fixed using pre-shock information; classifying a recession using output affected by the shock changes the conditioning event.

A single instrumental-variable ratio addresses a narrower object. In one state let observed cumulative variables be $Y,G$, and let instrument $Z$ satisfy the declared homogeneous moment model
\[
 \E[Z(Y-mG)\mid s]=0.
 \tag{1}
\]
Set $a_s=\E[ZY\mid s]$ and $b_s=\E[ZG\mid s]$. If $b_s\ne0$, the unique moment coefficient is $m_s=a_s/b_s$. The exclusion in (1) is not implied by calling $Z$ a fiscal news shock. It also does not identify a nonlinear response surface without additional restrictions.
\section{A confidence set that does not divide by relevance}
For a finite-sample benchmark observe $n$ independent episodes in the same state, with $|Z_i|\le1$ and $|Y_i|,|G_i|\le M$. Episodes are separated sufficiently to make the assumed independence meaningful; overlapping horizon sums from consecutive dates are not automatically independent.

Let $\widehat a=n^{-1}\sum_iZ_iY_i$, $\widehat b=n^{-1}\sum_iZ_iG_i$, and
$t_n=\sqrt{2\log(2/\alpha)/n}$. Define
\[
 \mathcal C_n=\{m\in\R:
 |\widehat a-m\widehat b|\le M(1+|m|)t_n\}.
 \tag{2}
\]
\begin{theorem}
For every true finite $m$ satisfying (1), the confidence set (2) covers it with probability at least $1-\alpha$, uniformly over the bounded independent law class. No lower bound on $|b_s|$ is required.
\end{theorem}
\begin{proof}
At the true coefficient, $X_i=Z_i(Y_i-mG_i)$ has mean zero and absolute bound $M(1+|m|)$. The bounded-mean exponential inequality gives
$\Prob(|n^{-1}\sum_iX_i|>M(1+|m|)t_n)\le\alpha$.
On the complementary event the defining inequality in (2) holds at that coefficient. Since the bound has the same probability for every admissible law, the coverage is uniform. Inverting one valid test at each candidate does not require a union bound over all real candidates to cover the true one.
\end{proof}
When the first stage is weak, (2) may be unbounded or contain two rays. That is compatible with valid uncertainty. Imposing an arbitrary finite plotting range and then describing the clipped interval as the full confidence set would hide the weak-identification problem.

A second construction starts from simultaneous intervals for $a_s,b_s$ and projects all ratios $a/b$ with $b\ne0$. If the first-stage interval is bounded away from zero, the ratio extrema occur among its four corner ratios. If it crosses zero, division of interval endpoints is generally incorrect: the set can be unbounded or disconnected. The projection is an uncertainty construction for the moment ratio, not proof that it equals a large spending intervention multiplier.
\section{What two randomized nonzero doses identify}
Consider the restricted state-specific response model
\[
 \E[G^\Sigma(a)\mid s]=a,\qquad
 \E[Y^\Sigma(a)\mid s]=\beta_s a+\gamma_s a^2
 \tag{3}
\]
for doses in a known interval containing zero. Purchase flows and output are in comparable units, and all compared doses use the same financing and monetary rule. Negative doses represent purchase reductions under that rule.

For $a\ne0$ the average multiplier is
$m_s(a)=\beta_s+\gamma_s a$.
The derivative multiplier at zero is $\beta_s$, whereas the derivative ratio at nonzero $a$ is $\beta_s+2\gamma_sa$. Positive and negative shocks of equal magnitude differ by $2\gamma_sa$ in their average multipliers.

\begin{proposition}
If randomized doses $0,a_1,a_2$ have positive assignment probabilities, with $a_1a_2\ne0$ and $a_1\ne a_2$, the model identifies
\[
 \gamma_s=
 \frac{\mu_s(a_2)/a_2-\mu_s(a_1)/a_1}{a_2-a_1},
 \qquad
 \beta_s=\mu_s(a_1)/a_1-\gamma_sa_1,
 \tag{4}
\]
where $\mu_s(a)=\E[Y^\Sigma(a)\mid s]$.
\end{proposition}
\begin{proof}
Randomization identifies each mean response relative to dose zero. Dividing (3) at the two nonzero doses yields two linear equations
$\mu_s(a_j)/a_j=\beta_s+\gamma_sa_j$. Their determinant is $a_2-a_1\ne0$. Subtraction and substitution give (4).
\end{proof}
With only dose zero and dose one, a response of $1.2$ is compatible with $(\beta,\gamma)=(1.2,0)$ or $(0.8,0.4)$. They agree on all tested outcomes but predict multipliers $1.2$ and $0.4$ respectively for dose minus one. One can construct identical outcome noise at the observed doses, making the full observed laws coincide. The missing negative-shock conclusion is therefore an extrapolation, not a small-sample issue.

Even two nonzero doses identify (4) only under the quadratic restriction. Adding a function proportional to $a(a-a_1)(a-a_2)$ leaves the tested means unchanged while altering other doses. A sufficiently small coefficient preserves any strict local outcome restrictions. Finite support does not identify an unrestricted smooth response surface everywhere.
\section{Spending composition can mimic state dependence}
Suppose investment purchases have constant multiplier two and current consumption purchases have constant multiplier one half. In each state an instrument induces a positive investment purchase response $g_I$ and consumption response $g_C$. With additive output effects its instrument ratio is
\[
 m^{IV}=\frac{2g_I+\tfrac12g_C}{g_I+g_C}.
\]
If the investment share is $0.8$ in recessions and $0.2$ in normal periods, the ratios are $1.7$ and $0.8$. Neither component multiplier varies by state. The difference comes entirely from instrument composition.

The same logic applies to financing, anticipation, and monetary accommodation. Comparing ratios across states requires either the same intervention mix or an explicit decomposition identifying how each component changes. If an instrument induces components with opposite signs, the weights need not be convex and the ratio can lie outside the component-multiplier range.

\section{Remaining empirical and dynamic work}
The moment confidence set, quadratic identification, and composition example are complete within their declared classes. They do not validate a fiscal instrument, identify the full response path, establish independence of historical episodes, or transport crisis-dose responses to ordinary times. Continuous states require support and smoothness assumptions, while overlapping outcomes need dependence-aware inference.

The broad catalogue problem remains unresolved. In particular, a narrow confidence interval for one well-identified local ratio is not evidence that all shock sizes, signs, states, financing arrangements, and purchase categories share its value. These are elementary restricted results rather than claims of a new general fiscal multiplier theory.
\begin{thebibliography}{9}
\bibitem{ars} A. Auclert, M. Rognlie, and L. Straub (2025).
\emph{Fiscal and Monetary Policy with Heterogeneous Agents}, section 5.
\url{https://web.stanford.edu/~aauclert/annual_review_hank.pdf}.
The primary review discusses nonlinear size and sign questions; equations (1)--(4) are the explicit models used here.
\end{thebibliography}

\end{document}
